\subsection{ANTI-INVARIANT ELEMENTS}

We retain the notations of the preceding number. Denote by \(\varepsilon(w)\) the determinant of the element \(w \in W\) . Thus

\[
\varepsilon (w) = (- 1) ^ {l (w)},
\]

the length \(l(w)\) being taken relative to the family of reflections \(s_{\alpha_i}\) .

DEFINITION 2. An element \(x \in \mathrm{A}[\mathrm{P}]\) is said to be anti-invariant under W if

\[
w (x) = (- 1) ^ {l (w)}. x
\]

for all \(w \in W\) .

The anti-invariant elements of A{[}P{]} form an A-submodule of A{[}P{]}. For any \(x \in A[P]\) , put

\[
J (x) = \sum_ {w \in W} \varepsilon (w). w (x).\tag{1}
\]

For \(x \in \mathrm{A}[\mathrm{P}]\) and \(w \in \mathrm{W}\) , we have

\[
w (\mathrm{J} (x)) = \sum_ {v \in \mathrm{W}} \varepsilon (v). w v (x) = \varepsilon (w) \sum_ {v \in \mathrm{W}} \varepsilon (v). v (x) = \varepsilon (w). \mathrm{J} (x)
\]

and \(J(x)\) is anti-invariant. On the other hand, let \(q = \text{Card}(W)\) . For any anti-invariant element x of A{[}P{]}, we have \(J(x) = q.x\) . It follows that, if q is invertible in A, the map \(q^{-1}J\) is a projection from A{[}P{]} onto the submodule of anti-invariant elements.

Let \(\overline{\omega}_{1},\ldots,\overline{\omega}_{l}\) be the fundamental weights corresponding to the chamber C. The elements of \(P\cap\overline{C}\) (resp. \(P\cap C\)) are the weights of the form \(n_{1}\overline{\omega}_{1}+\cdots+n_{l}\overline{\omega}_{l}\) with \(n_{i}\geqslant0\) (resp. \(n_{i}>0\) ) for \(1\leqslant i\leqslant l\) (§ 1, no. 10). On the other hand,

\[
\rho = \overline {{\omega}} _ {1} + \dots + \overline {{\omega}} _ {l}
\]

is half the sum of the positive roots (loc. cit.) so the elements of \(\mathrm{P} \cap \mathrm{C}\) are the weights of the form \(\rho + p\) with \(p \in \mathrm{P} \cap \overline{\mathrm{C}}\) . Finally, if \(p \in \mathrm{P} \cap \mathrm{C}\) , then \(w(p) < p\) for all \(w \neq 1\) ( \(\S 1\) , no. 6, Cor. to Prop. 18) and \(e^p\) is thus the unique maximal term of \(\mathrm{J}(e^p)\) .

PROPOSITION 1. If 2 is not a zero divisor in A, the elements \(\mathrm{J}(e^{p})\) for \(p \in P \cap C\) form a basis of the module of anti-invariant elements of A{[}P{]}.

The weights \(w(p)\) for \(w \in \mathbf{W}\) and \(p \in \mathrm{P} \cap \mathrm{C}\) are pairwise distinct. It follows that the \(\mathrm{J}(e^p)\) for \(p \in \mathrm{P} \cap \mathrm{C}\) are linearly independent.

On the other hand, let \(x = \sum_{p} x_{p} e^{p}\) be an anti-invariant element of \(A[P]\) . If \(p_{0}\) belongs to a wall, it is invariant under a reflection \(s \in W\) and

\[
x = \sum_ {p} x _ {p} e ^ {p} = - s (x) = - \sum_ {p} x _ {p} e ^ {s (p)}.
\]

It follows that \(2x_{p_{0}} = 0\) , so \(x_{p_{0}} = 0\) . Since every element that does not belong to any wall can be written uniquely in the form \(w(p)\) with \(w \in W\) and \(p \in P \cap C\) , we thus have

\[
x = \sum_ {p \in P \cap C} \sum_ {w \in W} x _ {w (p)} e ^ {w (p)}.\tag{2}
\]

Since \(w(x) = \sum_{p} x_{p} e^{w(p)} = \varepsilon(w) \sum_{p} x_{p} e^{p}, x_{w(p)} = \varepsilon(w) x_{p}\) and we deduce from (2) that

\[
x = \sum_ {p \in P \cap C} x _ {p} J (e ^ {p}),
\]

which completes the proof.

Consider now the element \(d\) of the algebra \(\mathrm{A}[\frac{1}{2}\mathrm{P}]\) defined by

\[
\begin{array}{l} d = \prod_ {\alpha \in R, \alpha > 0} (e ^ {\alpha / 2} - e ^ {- \alpha / 2}) \\ \quad = e ^ {\rho}. \prod_ {\alpha \in R, \alpha > 0} (1 - e ^ {- \alpha}) \\ \quad = e ^ {- \rho}. \prod_ {\alpha \in R, \alpha > 0} (e ^ {\alpha} - 1). \end{array}\tag{3}
\]

Since \(\rho \in \mathrm{P}\) , \(d \in \mathrm{A}[\mathrm{P}]\) .

PROPOSITION 2. (i) The element \(d\) defined by (3) is an anti-invariant element of A{[}P{]}; its unique maximal term (no. 2, Def. 1) is \(e^{\rho}\) and \(d = \mathrm{J}(e^{\rho})\) .

\begin{enumerate}
\def\labelenumi{(\roman{enumi})}
\setcounter{enumi}{1}
\item
  For any \(p \in \mathbb{P}\) , the element \(\mathrm{J}(e^p)\) is divisible uniquely by \(d\) and the quotient \(\mathrm{J}(e^p)/d\) is an element of \(\mathrm{A}[\mathrm{P}]\) invariant under \(\mathrm{W}\) .
\item
  If 2 is not a zero divisor in A, multiplication by d is a bijection from the set of elements of A{[}P{]} invariant under W to the set of anti-invariant elements of A{[}P{]}.
\end{enumerate}

We know that, for \(1 \leqslant i \leqslant l\) , the reflection \(s_i = s_{\alpha_i}\) leaves stable the set of positive roots other than \(\alpha_i\) and that \(s_i(\alpha_i) = -\alpha_i\) (§ 1, no. 6. Cor. 1 of Prop. 17). Hence,

\[
\begin{array}{l} s _ {i} (d) = (e ^ {- \alpha_ {i} / 2} - e ^ {\alpha_ {i} / 2}). \prod_ {\alpha \in R, \alpha > 0, \alpha \neq \alpha_ {i}} (e ^ {\alpha / 2} - e ^ {- \alpha / 2}) \\ = - d = \varepsilon (s _ {i}). d. \end{array}
\]

Since the \(s_{i}\) generate W, this proves the first assertion in (i). The second assertion in (i) follows immediately from (3) and Lemma 2, noting that 1 is the unique maximal term of \(1 - e^{-\alpha}\) for \(\alpha \in R\) , \(\alpha > 0\) .

Assume now that A = Z. By Prop. 1,

\[
d = \prod_ {p \in \mathrm{P} \cap \mathrm{C}} c _ {p} \mathrm{J} (e ^ {p}) \text {with} c _ {p} \in \mathbf {Z}.\tag{4}
\]

On the other hand, it is clear that

\[
d = e ^ {\rho} + \sum_ {q <   \rho} c _ {q} ^ {\prime} e ^ {q}.\tag{5}
\]

If \(p \in \mathrm{P} \cap \mathrm{C}\) with \(p \neq \rho\) , then \(p > \rho\) and the coefficient of \(e^p\) in \(d\) is zero by (5). Thus, \(c_p = 0\) . Moreover, comparison of the coefficients of \(e^\rho\) in (4) and (5) shows that \(c_\rho = 1\) and hence that \(d = \mathrm{J}(e^\rho)\) .

We continue to assume that A = Z. Let \(p \in P\) , \(\alpha \in R\) and M be a system of representatives of the right cosets of W with respect to the subgroup \(\{1, s_{\alpha}\}\) . Then,

\[
\mathrm{J} (e ^ {p}) = \sum_ {w \in \mathrm{M}} \varepsilon (w) e ^ {w (p)} + \sum_ {w \in \mathrm{M}} \varepsilon (s _ {\alpha} w) e ^ {s _ {\alpha} w (p)}.
\]

Now \(s_{\alpha}w(p) = w(p) - \langle \alpha^{\prime},w(p)\rangle \alpha = w(p) + n_w\alpha\) , with \(n_w\in \mathbf{Z}\) . Thus,

\[
J (e ^ {p}) = \sum_ {w \in M} \varepsilon (w) e ^ {w (p)} \left(1 - e ^ {n _ {w} \alpha}\right).
\]

If \(n_w \geqslant 0\) , it is clear that \(1 - e^{n_w\alpha}\) is divisible by \(1 - e^\alpha\) and this is also true when \(n_w < 0\) since \(1 - e^{n_w\alpha} = -e^{n_w\alpha}(1 - e^{-n_w\alpha})\) . Hence, \(J(e^p)\) is divisible by \(1 - e^\alpha\) in \(\mathbf{Z}[P]\) .

By Lemma 1, \(\mathbf{Z}[\mathrm{P}]\) is factorial and the elements \(1 - e^{\alpha}\) for \(\alpha \in \mathbb{R}\) and \(\alpha > 0\) are mutually prime. It follows that \(\mathrm{J}(e^{p})\) is divisible in \(\mathbf{Z}[\mathrm{P}]\) by the product \(\prod_{\alpha > 0}(1 - e^{\alpha})\) , and hence also by \(d = e^{-\rho}\prod_{\alpha > 0}(1 - e^{\alpha})\) .

Returning to the general case, by extension of scalars from Z to A, we deduce from the above that \(d = \mathrm{J}(e^{\rho})\) and that every element \(\mathrm{J}(e^{p})\) is divisible by d.~Since \(e^{\rho}\) is the unique maximal term of d, the Remark of no. 2 shows that there exists a unique element \(y \in A[P]\) such that \(\mathrm{J}(e^{p}) = dy\) and it follows immediately that y is invariant under W, and hence that d and \(\mathrm{J}(e^{p})\) are anti-invariant. This proves (i) and (ii).

Finally, if 2 is not a zero divisor in A, the Remark of no. 2 and Prop. 1 imply (iii).

Remarks. 1) If 2 is not a zero divisor in A, it is easy to check that d is the unique anti-invariant element of A{[}P{]} with \(e^{\rho}\) as its maximal term.

\begin{enumerate}
\def\labelenumi{\arabic{enumi})}
\setcounter{enumi}{1}
\tightlist
\item
  Lemma 2 of no. 2 shows that the unique maximal term of the quotient \(\mathrm{J}(e^p) / d\) (for \(p\in \mathrm{P}\cap \mathrm{C}\) ) is \(e^{p - \rho}\) .
\end{enumerate}

